\chapter{Methodology}
\label{ch:methodology}

This chapter presents the technical foundations of the \acrshort{piec} Navigation Framework. We begin with an overview of the system architecture in Section~\ref{sec:meth_architecture}, then detail each of the five key components: the \acrshort{pinn} surrogate model (Section~\ref{sec:meth_pinn}), multi-objective \acrshort{nsga} path optimizer (Section~\ref{sec:meth_optimizer}), \acrshort{ukf} localization (Section~\ref{sec:meth_ukf}), adaptive \acrshort{dwa} controller (Section~\ref{sec:meth_dwa}), and system integration (Section~\ref{sec:meth_integration}).

\section{System Architecture}
\label{sec:meth_architecture}

The \acrshort{piec} framework consists of five integrated components operating in a hierarchical control architecture (see Figure~\ref{fig:system_architecture}):

\begin{enumerate}
    \item \textbf{Perception Layer:} Sensor inputs (LiDAR, IMU, wheel encoders, magnetometer)
    \item \textbf{State Estimation:} \acrshort{ukf} localization providing 5D state estimate $\vect{x} = [x, y, \theta, v, \omega]\transpose$
    \item \textbf{Global Planning:} \acrshort{nsga} multi-objective path optimizer
    \item \textbf{Energy Prediction:} \acrshort{pinn} surrogate model service
    \item \textbf{Local Control:} Adaptive \acrshort{dwa} obstacle avoidance
\end{enumerate}

\subsection{Information Flow}

The system operates at multiple temporal frequencies:
\begin{itemize}
    \item \textbf{Control loop:} 20 Hz (50ms cycle) - \acrshort{dwa} trajectory execution
    \item \textbf{State estimation:} 20 Hz (50ms cycle) - \acrshort{ukf} updates
    \item \textbf{Global planning:} 1 Hz (1s cycle) - \acrshort{nsga} replanning
    \item \textbf{PINN queries:} On-demand (typically 10-15 calls per planning cycle)
\end{itemize}

% TODO: Add Figure 3.1 - System architecture diagram showing:
% - Sensor inputs → UKF → State estimate
% - Goal + Map → NSGA Optimizer → Global path
% - PINN Service (queried by optimizer)
% - Global path + Local sensors → DWA → Velocity commands
% - Robot platform
% Reference: Will be created as figures/system_architecture.pdf

\subsection{Component Interactions}

\paragraph{UKF → Optimizer} State estimate $\vect{x}$ and covariance $\mat{P}$ provide current robot pose and uncertainty for path planning initialization.

\paragraph{Optimizer ↔ PINN} For each candidate path in \acrshort{nsga}, the optimizer queries the \acrshort{pinn} service to predict energy consumption and stability, which become objectives in the multi-objective optimization.

\paragraph{Optimizer → DWA} The selected Pareto-optimal path from \acrshort{nsga} serves as the global reference for local trajectory tracking.

\paragraph{DWA → Robot} The \acrshort{dwa} controller outputs velocity commands $(v, \omega)$ at 20 Hz for real-time obstacle avoidance.

\section{PINN Surrogate Model}
\label{sec:meth_pinn}

The \acrshort{pinn} surrogate model predicts mobile robot energy consumption and stability by learning from data while respecting physical constraints. This section derives the complete formulation.

% Code Reference: src/piec_pinn_surrogate/piec_pinn_surrogate/pinn_model.py (228 lines)

\subsection{Input-Output Specification}

\paragraph{Input Features}
The network accepts a 10-dimensional input vector $\vect{u} \in \Real^{10}$:
\begin{equation}
\vect{u} = \begin{bmatrix}
x \\ y \\ \theta \\ v \\ \omega \\ \text{slope} \\ \text{roughness} \\ \rho_{\text{obs}} \\ d_{\text{clear}} \\ \text{terrain}
\end{bmatrix}
\label{eq:pinn_input}
\end{equation}

where:
\begin{itemize}
    \item $x, y \in \Real$: Robot position (m)
    \item $\theta \in [-\pi, \pi]$: Heading angle (rad)
    \item $v \in [0, v_{\max}]$: Linear velocity (m/s)
    \item $\omega \in [-\omega_{\max}, \omega_{\max}]$: Angular velocity (rad/s)
    \item $\text{slope} \in [-1, 1]$: Terrain slope (normalized)
    \item $\text{roughness} \in [0, 1]$: Surface roughness
    \item $\rho_{\text{obs}} \in [0, 1]$: Obstacle density from laser scan
    \item $d_{\text{clear}} \in [0, d_{\max}]$: Minimum clearance to obstacles (m)
    \item $\text{terrain} \in \{0, 1, 2, \ldots\}$: Discrete terrain type
\end{itemize}

\paragraph{Output Features}
The network outputs a 2-dimensional vector $\vect{y} \in \Real^2$:
\begin{equation}
\vect{y} = \begin{bmatrix}
E \\ S
\end{bmatrix}
\label{eq:pinn_output}
\end{equation}

where:
\begin{itemize}
    \item $E \geq 0$: Energy consumption rate (J/m), constrained via ReLU activation
    \item $S \in [0, 1]$: Stability metric, constrained via sigmoid activation
\end{itemize}

\subsection{Network Architecture}

The \acrshort{pinn} uses a fully connected neural network with the following architecture:

\begin{align}
\vect{h}_1 &= \tanh(\mat{W}_1 \vect{u} + \vect{b}_1) \quad &\text{(Input → 128 units)} \label{eq:layer1}\\
\vect{h}'_1 &= \text{Dropout}(\vect{h}_1, p=0.1) \quad &\text{(Regularization)} \label{eq:dropout1}\\
\vect{h}_2 &= \tanh(\mat{W}_2 \vect{h}'_1 + \vect{b}_2) \quad &\text{(128 → 128 units)} \label{eq:layer2}\\
\vect{h}'_2 &= \text{Dropout}(\vect{h}_2, p=0.1) \quad &\text{(Regularization)} \label{eq:dropout2}\\
\vect{h}_3 &= \tanh(\mat{W}_3 \vect{h}'_2 + \vect{b}_3) \quad &\text{(128 → 64 units)} \label{eq:layer3}\\
\vect{y} &= \mat{W}_4 \vect{h}_3 + \vect{b}_4 \quad &\text{(64 → 2 units)} \label{eq:output_layer}
\end{align}

\textbf{Design Rationale:}
\begin{itemize}
    \item \textbf{Tanh activation:} Smooth, bounded nonlinearity suitable for physics-informed learning
    \item \textbf{Dropout (p=0.1):} Prevents overfitting on training data
    \item \textbf{Gradual dimension reduction:} 128→128→64→2 preserves information while reducing parameters
    \item \textbf{No activation on output:} Allows unconstrained prediction before physics constraints
\end{itemize}

\subsection{Physics-Based Energy Decomposition}

The energy model decomposes total energy into six physical components:

\subsubsection{Kinetic Energy}
\begin{equation}
E_{\text{kinetic}} = \frac{1}{2} m v^2
\label{eq:kinetic_energy}
\end{equation}
Represents translational kinetic energy. For unit mass ($m=1$), this simplifies to $\frac{1}{2}v^2$.

\subsubsection{Turning Energy}
\begin{equation}
E_{\text{turning}} = \alpha_{\omega} |\omega|
\label{eq:turning_energy}
\end{equation}
where $\alpha_{\omega} = 0.2$ accounts for energy required for angular motion (wheel differential drive, friction during turns).

\subsubsection{Slope Energy}
\begin{equation}
E_{\text{slope}} = \alpha_s |\theta_{\text{slope}}| \cdot v
\label{eq:slope_energy}
\end{equation}
where $\alpha_s = 2.0$ and $\theta_{\text{slope}}$ is terrain inclination. Positive slopes require additional energy to overcome gravity: $F_{\text{gravity}} = mg\sin(\theta)$.

\subsubsection{Friction Energy}
\begin{equation}
E_{\text{friction}} = \alpha_r \cdot \text{roughness} \cdot v
\label{eq:friction_energy}
\end{equation}
where $\alpha_r = 0.1$. Rolling resistance depends on surface properties and velocity.

\subsubsection{Obstacle Energy}
\begin{equation}
E_{\text{obstacle}} = \alpha_{\text{obs}} \cdot \rho_{\text{obs}} \cdot v
\label{eq:obstacle_energy}
\end{equation}
where $\alpha_{\text{obs}} = 0.3$. Higher obstacle density increases path complexity, requiring more frequent velocity adjustments and turns.

\subsubsection{Clearance Energy}
\begin{equation}
E_{\text{clearance}} = \begin{cases}
-\alpha_c \cdot \frac{\text{ReLU}(d_{\text{safe}} - d_{\text{clear}})}{d_{\text{safe}}} & \text{if } d_{\text{clear}} < d_{\text{safe}} \\
0 & \text{otherwise}
\end{cases}
\label{eq:clearance_energy}
\end{equation}
where $\alpha_c = 0.1$ and $d_{\text{safe}} = 2.0$ m. This penalty term discourages paths close to obstacles (negative contribution, increasing total energy).

\subsubsection{Combined Physics Energy}
\begin{equation}
E_{\text{physics}} = E_{\text{kinetic}} + E_{\text{turning}} + E_{\text{slope}} + E_{\text{friction}} + E_{\text{obstacle}} + E_{\text{clearance}}
\label{eq:total_physics_energy}
\end{equation}

\subsection{Adaptive Energy Blending}

The final energy prediction adaptively blends neural network output $E_{\text{NN}}$ with physics-based energy $E_{\text{physics}}$:

\begin{equation}
w_{\text{physics}} = 0.3 + 0.3 \cdot \rho_{\text{obs}}
\label{eq:blend_weight}
\end{equation}

\begin{equation}
E_{\text{final}} = (1 - w_{\text{physics}}) E_{\text{NN}} + w_{\text{physics}} E_{\text{physics}}
\label{eq:energy_blending}
\end{equation}

\textbf{Rationale:} In high obstacle density regions ($\rho_{\text{obs}} \rightarrow 1$), physics models are more reliable (weight $\rightarrow 0.6$). In open space ($\rho_{\text{obs}} \rightarrow 0$), learned model dominates (weight $\rightarrow 0.3$).

\subsection{Stability Model}

The stability metric $S \in [0, 1]$ quantifies trajectory confidence:

\begin{equation}
S_{\text{reduction}} = 0.2|v| + 0.3|\theta_{\text{slope}}| + 0.4\rho_{\text{obs}} + 0.2\left(1 - \sigma(d_{\text{clear}} - 1)\right)
\label{eq:stability_reduction}
\end{equation}

where $\sigma(x) = \frac{1}{1 + e^{-x}}$ is the sigmoid function.

\begin{equation}
S_{\text{adjusted}} = \text{clamp}\left( S_{\text{NN}} (1 - S_{\text{reduction}}), 0.1, 1.0 \right)
\label{eq:stability_adjusted}
\end{equation}

\textbf{Interpretation:} Stability decreases with high velocity, steep slopes, dense obstacles, and small clearances. The base neural network prediction $S_{\text{NN}}$ is modulated by physics-based reduction factors.

\subsection{Physics-Informed Loss Function}

The total training loss combines data-fitting and physics constraints:

\begin{equation}
\mathcal{L}_{\text{total}} = \mathcal{L}_{\text{data}} + \lambda_{\text{physics}} \mathcal{L}_{\text{physics}}
\label{eq:total_loss}
\end{equation}

where $\lambda_{\text{physics}} = 0.15$ is the physics loss weight.

\subsubsection{Data Loss}
Standard mean squared error on labeled training data:
\begin{equation}
\mathcal{L}_{\text{data}} = \frac{1}{N} \sum_{i=1}^N \left\| \vect{y}_i - \hat{\vect{y}}_i \right\|^2
\label{eq:data_loss}
\end{equation}

\subsubsection{Physics Loss Components}

The physics loss enforces six constraints:

\paragraph{1. Non-Negativity Constraint}
Energy must be non-negative:
\begin{equation}
\mathcal{L}_{\text{non-neg}} = 0.1 \cdot \mathbb{E}\left[ \text{ReLU}(-E) \right]
\label{eq:loss_nonneg}
\end{equation}

\paragraph{2. Bounds Constraint}
Stability must be in $[0, 1]$:
\begin{equation}
\mathcal{L}_{\text{bounds}} = 0.1 \cdot \mathbb{E}\left[ \text{ReLU}(S - 1) + \text{ReLU}(-S) \right]
\label{eq:loss_bounds}
\end{equation}

\paragraph{3. Energy Conservation}
Predicted energy should match physics-based expectation:
\begin{equation}
\mathcal{L}_{\text{conservation}} = 0.01 \cdot \mathbb{E}\left[ (E_{\text{pred}} - E_{\text{physics}})^2 \right]
\label{eq:loss_conservation}
\end{equation}

\paragraph{4. Stability Physics}
Predicted stability should align with physics-based stability:
\begin{equation}
\mathcal{L}_{\text{stability}} = 0.01 \cdot \mathbb{E}\left[ (S_{\text{pred}} - S_{\text{expected}})^2 \right]
\label{eq:loss_stability}
\end{equation}

\paragraph{5. Clearance Penalty}
Discourage predictions favoring low clearances:
\begin{equation}
\mathcal{L}_{\text{clearance}} = 0.05 \cdot \mathbb{E}\left[ \text{ReLU}(0.5 - d_{\text{clear}}) \right]
\label{eq:loss_clearance}
\end{equation}

\paragraph{6. Obstacle Smoothness}
Regularize obstacle density influence:
\begin{equation}
\mathcal{L}_{\text{smoothness}} = 0.01 \cdot \mathbb{E}\left[ |\rho_{\text{obs}}| \right]
\label{eq:loss_smoothness}
\end{equation}

\paragraph{Combined Physics Loss}
\begin{equation}
\mathcal{L}_{\text{physics}} = \mathcal{L}_{\text{non-neg}} + \mathcal{L}_{\text{bounds}} + \mathcal{L}_{\text{conservation}} + \mathcal{L}_{\text{stability}} + \mathcal{L}_{\text{clearance}} + \mathcal{L}_{\text{smoothness}}
\label{eq:physics_loss_combined}
\end{equation}

\subsection{Training Procedure}

\textbf{Optimizer:} Adam with learning rate $\eta = 10^{-3}$, weight decay $\lambda_{\text{wd}} = 10^{-5}$

\textbf{Scheduler:} ReduceLROnPlateau with patience=10, factor=0.5, min\_lr=$10^{-6}$

\textbf{Batch size:} 64

\textbf{Epochs:} 200 (with early stopping, patience=20)

\textbf{Hardware:} NVIDIA GPU (CUDA acceleration)

\textbf{Training time:} Approximately 15-20 minutes for 10,000 samples

% TODO: Add Figure 3.2 - PINN network architecture diagram
% Show: Input (10D) → FC1 (128) → FC2 (128) → FC3 (64) → Output (2D)
% Include activation functions and dropout layers
% Reference: Will be created as figures/pinn_architecture.pdf

% TODO: Add Figure 3.3 - Energy component decomposition
% Bar chart showing contribution of each of 6 energy components
% Reference: Will be created as figures/energy_decomposition.pdf


\section{Multi-Objective NSGA-II Path Optimizer}
\label{sec:meth_optimizer}

The path optimizer uses \acrfull{nsga} to simultaneously optimize seven competing objectives, producing Pareto-optimal path sets from which mission-specific solutions can be selected.

% Code Reference: src/piec_path_optimizer/piec_path_optimizer/
% - complete_path_optimizer.py (2348 lines)
% - objectives.py, nsga2.py, uncertainty_aware.py

\subsection{Problem Formulation}

\paragraph{Path Representation}
A path $\mathcal{P}$ is represented as a sequence of waypoints:
\begin{equation}
\mathcal{P} = \{(x_0, y_0), (x_1, y_1), \ldots, (x_n, y_n)\}
\label{eq:path_representation}
\end{equation}
where $(x_0, y_0)$ is the start position, $(x_n, y_n)$ is the goal, and intermediate waypoints $n \leq 8$ are optimized.

\paragraph{Multi-Objective Optimization Problem}
\begin{equation}
\begin{aligned}
\min_{\mathcal{P}} \quad & \vect{f}(\mathcal{P}) = [f_1, f_2, \ldots, f_7]\transpose \\
\text{s.t.} \quad & \mathcal{P} \text{ is collision-free} \\
& d_{\text{clear}}(\mathcal{P}) \geq d_{\min} \\
& \mathcal{P}_0 = \text{start}, \quad \mathcal{P}_n = \text{goal}
\end{aligned}
\label{eq:moo_problem}
\end{equation}

\subsection{Objective Functions}

\subsubsection{Objective 1: Path Length}
\begin{equation}
f_1 = L(\mathcal{P}) = \sum_{i=0}^{n-1} \sqrt{(x_{i+1} - x_i)^2 + (y_{i+1} - y_i)^2}
\label{eq:obj_length}
\end{equation}

\textbf{Weight:} $w_1 = 0.12$. Minimizing path length reduces travel time and baseline energy consumption.

\subsubsection{Objective 2: Average Curvature}
\begin{equation}
f_2 = \frac{1}{n-2} \sum_{i=1}^{n-2} \arccos\left( \frac{\vect{v}_i \cdot \vect{v}_{i+1}}{\|\vect{v}_i\| \|\vect{v}_{i+1}\|} \right)
\label{eq:obj_curvature}
\end{equation}

where $\vect{v}_i = (x_{i+1} - x_i, y_{i+1} - y_i)$ is the path segment vector.

\textbf{Weight:} $w_2 = 0.08$. Smooth paths with low curvature reduce angular accelerations and improve tracking accuracy.

\subsubsection{Objective 3: Obstacle Cost}
\begin{equation}
f_3 = C_{\text{obs}} - 0.3 \cdot \text{bonus}_{\text{free}}
\label{eq:obj_obstacle_main}
\end{equation}

where the obstacle cost is:
\begin{equation}
C_{\text{obs}} = \sum_{i=0}^{n} w_{\text{penalty}} \cdot \max(0, d_{\min} - d_{\text{clear}}(x_i, y_i))^2
\label{eq:obstacle_cost}
\end{equation}

with $w_{\text{penalty}} = 6.0$ and $d_{\min} = 0.4$ m.

The free-space bonus rewards paths through open regions:
\begin{equation}
\text{bonus}_{\text{free}} = \frac{1}{n} \sum_{i=0}^{n-1} \sum_{j=0}^{N_{\text{samples}}} \mathbb{I}_{\text{free}}(x_j, y_j)
\label{eq:free_space_bonus}
\end{equation}

where $\mathbb{I}_{\text{free}} = 1$ if free-space grid value $> 0.7$, else 0.

\textbf{Weight:} $w_3 = 0.35$ (highest weight). Safety is paramount in navigation.

\subsubsection{Objective 4: Localization Uncertainty}
\begin{equation}
f_4 = \frac{1}{n} \sum_{i=0}^{n} U(x_i, y_i)
\label{eq:obj_uncertainty}
\end{equation}

where the uncertainty at each point is:
\begin{equation}
U(x, y) = \text{trace}(\mat{\Sigma}_{\text{pos}}(x, y)) + 0.5 \sigma_{\text{yaw}}^2(x, y)
\label{eq:uncertainty_metric}
\end{equation}

$\mat{\Sigma}_{\text{pos}}$ is the position covariance from the uncertainty map, and $\sigma_{\text{yaw}}^2$ is yaw variance.

\textbf{Weight:} $w_4 = 0.08$. Paths through high-uncertainty regions degrade localization quality.

\subsubsection{Objective 5: Deviation from Straight Line}
\begin{equation}
f_5 = \frac{|L(\mathcal{P}) - L_{\text{straight}}|}{\max(L_{\text{straight}}, 0.1)}
\label{eq:obj_deviation}
\end{equation}

where $L_{\text{straight}} = \|(x_n, y_n) - (x_0, y_0)\|$ is the Euclidean distance.

\textbf{Weight:} $w_5 = 0.12$. Encourages paths closer to direct routes when feasible.

\subsubsection{Objective 6: PINN-Predicted Energy}
\begin{equation}
f_6 = E_{\text{PINN}}(\mathcal{P}) = \sum_{i=0}^{n-1} E_{\text{PINN}}(x_i, y_i, \theta_i, v_i, \omega_i, \ldots)
\label{eq:obj_pinn_energy}
\end{equation}

The \acrshort{pinn} service is queried for each path segment with estimated velocities and terrain properties.

\textbf{Weight:} $w_6 = 0.15$ (second highest). Energy efficiency is a primary objective.

\subsubsection{Objective 7: PINN-Predicted Stability}
\begin{equation}
f_7 = 1.0 - \frac{1}{n} \sum_{i=0}^{n-1} S_{\text{PINN}}(x_i, y_i, \ldots)
\label{eq:obj_pinn_stability}
\end{equation}

Stability values are in $[0, 1]$. We minimize $(1 - S)$ to maximize stability.

\textbf{Weight:} $w_7 = 0.10$. Stable paths reduce risk of localization failure.

\subsection{Composite Fitness Function}

For selection purposes (not the multi-objective optimization itself), we define a scalar fitness:
\begin{equation}
F(\mathcal{P}) = \sum_{j=1}^{7} w_j \cdot f_j(\mathcal{P})
\label{eq:composite_fitness}
\end{equation}

This is used only for reporting and initial population ranking; \acrshort{nsga} operates on the full objective vector.

\subsection{NSGA-II Algorithm}

\subsubsection{Fast Non-Dominated Sorting}

\paragraph{Dominance Relation}
Path $\mathcal{P}_a$ dominates $\mathcal{P}_b$ (denoted $\mathcal{P}_a \prec \mathcal{P}_b$) if:
\begin{equation}
\forall j: f_j(\mathcal{P}_a) \leq f_j(\mathcal{P}_b) \quad \land \quad \exists k: f_k(\mathcal{P}_a) < f_k(\mathcal{P}_b)
\label{eq:dominance}
\end{equation}

\paragraph{Sorting Algorithm}
For each solution $p$ in population $P$:
\begin{enumerate}
    \item Compute domination set $S_p = \{q \in P : p \prec q\}$
    \item Compute domination counter $n_p = |\{q \in P : q \prec p\}|$
\end{enumerate}

Front assignment:
\begin{equation}
\begin{aligned}
F_1 &= \{p \in P : n_p = 0\} \quad \text{(non-dominated solutions)} \\
F_{i+1} &= \{q \in S_p : p \in F_i, \, n_q' = 0\}
\end{aligned}
\label{eq:front_assignment}
\end{equation}

where $n_q' = n_q - 1$ after processing each front.

\textbf{Complexity:} $O(MN^2)$ for $M$ objectives and $N$ individuals.

\subsubsection{Crowding Distance}

For diversity preservation within each front $F_i$, crowding distance $d_c(p)$ is computed:

\paragraph{Algorithm}
\begin{enumerate}
    \item For each objective $m = 1, \ldots, M$:
    \begin{enumerate}
        \item Sort $F_i$ by $f_m$
        \item Set boundary distances: $d_c(\text{first}) = d_c(\text{last}) = \infty$
        \item For interior points $j = 2, \ldots, |F_i| - 1$:
        \begin{equation}
        d_c^{(m)}(j) = \frac{f_m(j+1) - f_m(j-1)}{f_m^{\max} - f_m^{\min}}
        \label{eq:crowding_distance_obj}
        \end{equation}
    \end{enumerate}
    \item Aggregate across objectives:
    \begin{equation}
    d_c(p) = \sum_{m=1}^{M} d_c^{(m)}(p)
    \label{eq:crowding_distance_total}
    \end{equation}
\end{enumerate}

\textbf{Interpretation:} Higher $d_c$ indicates the solution is in a sparse region of objective space, promoting diversity.

\subsubsection{Selection}

Tournament selection with crowded comparison operator $\prec_n$:

\begin{equation}
p \prec_n q \iff \begin{cases}
\text{rank}(p) < \text{rank}(q) & \text{(different fronts)} \\
d_c(p) > d_c(q) & \text{(same front, prefer diversity)}
\end{cases}
\label{eq:crowded_comparison}
\end{equation}

\subsubsection{Genetic Operators}

\paragraph{Crossover (Rate: 0.8)}
Two-point crossover exchanges waypoint subsequences between parent paths:
\begin{equation}
\begin{aligned}
\text{Child}_1 &= \text{Parent}_1[1:k_1] + \text{Parent}_2[k_1:k_2] + \text{Parent}_1[k_2:n] \\
\text{Child}_2 &= \text{Parent}_2[1:k_1] + \text{Parent}_1[k_1:k_2] + \text{Parent}_2[k_2:n]
\end{aligned}
\label{eq:crossover}
\end{equation}

where $k_1, k_2$ are randomly selected crossover points.

\paragraph{Mutation (Rate: 0.5)}
Uncertainty-aware mutation with free-space consideration (see Section~\ref{sec:uncertainty_mutation}).

\subsection{Uncertainty-Aware Mutation}
\label{sec:uncertainty_mutation}

Standard mutation perturbs waypoints randomly. Our enhancement biases mutation toward high free-space, low-uncertainty regions:

\paragraph{Mutation Algorithm}
For each waypoint $i$ (except start/goal):
\begin{enumerate}
    \item With probability $p_{\text{mut}} = 0.5$, mutate waypoint
    \item Evaluate current fitness:
    \begin{equation}
    \Phi_{\text{current}} = 0.6 \cdot C(x_i, y_i) + 0.4 \cdot F(x_i, y_i)
    \label{eq:mutation_fitness_current}
    \end{equation}
    where $C$ is clearance and $F$ is free-space score
    \item Sample $N_{\text{candidates}} = 5$ candidate positions:
    \begin{equation}
    (x'_j, y'_j) = (x_i, y_i) + r_j (\cos\phi_j, \sin\phi_j)
    \label{eq:candidate_sampling}
    \end{equation}
    where $r_j \sim \mathcal{U}(0, 0.6)$ m and $\phi_j \sim \mathcal{U}(0, 2\pi)$
    \item Select candidate with highest fitness:
    \begin{equation}
    (x_i, y_i) \leftarrow \arg\max_{j} \Phi_j
    \label{eq:best_candidate}
    \end{equation}
\end{enumerate}

\paragraph{Free-Space Score}
For a point $(x, y)$, evaluate $7 \times 7$ neighborhood:
\begin{equation}
F(x, y) = \frac{1}{49} \sum_{di=-3}^{3} \sum_{dj=-3}^{3} \mathbb{I}_{\text{free}}(\text{cell}_{x+di, y+dj})
\label{eq:free_space_score}
\end{equation}

This encourages mutation toward open spaces, improving path feasibility.

\subsection{Stuck Detection and Recovery}

Long-duration missions may encounter situations where the robot becomes stuck (local minima, deadlocks). We implement three stuck detection mechanisms:

\subsubsection{Velocity-Based Detection}

Track total path length traveled over time window $T_{\text{window}} = 15$ s:
\begin{equation}
M_{\text{total}}(t) = \sum_{i=0}^{N(t)-1} \|(x_{i+1}, y_{i+1}) - (x_i, y_i)\|
\label{eq:total_movement}
\end{equation}

where positions are sampled at 20 Hz.

\textbf{Stuck criterion:}
\begin{equation}
M_{\text{total}} < \tau_{\text{move}} = 0.25 \text{ m}
\label{eq:stuck_velocity}
\end{equation}

\subsubsection{Angular Oscillation Detection}

Compute radius of gyration over position history:
\begin{equation}
R_{\text{gyration}} = \frac{1}{N} \sum_{i=1}^{N} \sqrt{(x_i - \bar{x})^2 + (y_i - \bar{y})^2}
\label{eq:radius_gyration}
\end{equation}

where $(\bar{x}, \bar{y})$ is the centroid of recent positions.

\textbf{Stuck criterion:}
\begin{equation}
R_{\text{gyration}} < \tau_{\text{radius}} = 0.35 \text{ m}
\label{eq:stuck_angular}
\end{equation}

Indicates robot is rotating in place without making progress.

\subsubsection{Positional Stagnation}

\textbf{Combined stuck detection:}
\begin{equation}
\text{Stuck} \iff \begin{cases}
M_{\text{total}} < 0.25 \text{ m} \, \land \\
R_{\text{gyration}} < 0.35 \text{ m} \, \land \\
t - t_{\text{goal\_received}} > 12 \text{ s} \, \land \\
d_{\text{to\_goal}} > 0.35 \text{ m}
\end{cases}
\label{eq:stuck_combined}
\end{equation}

\textbf{Recovery action:} Trigger escape sequence (replan with exploration bias, increase free-space weight to 0.5).

\textbf{Cooldown:} 20 s between consecutive stuck detections to prevent rapid oscillation.

\subsection{Algorithm Parameters}

\begin{table}[h]
\centering
\caption{NSGA-II Path Optimizer Parameters}
\label{tab:nsga_params}
\begin{tabular}{lcc}
\toprule
\textbf{Parameter} & \textbf{Value} & \textbf{Unit} \\
\midrule
Population size & 40 & individuals \\
Generations & 12 & - \\
Crossover rate & 0.8 & - \\
Mutation rate & 0.5 & - \\
Max waypoints & 8 & points \\
Optimization timeout & 2.0 & s \\
Planning rate & 1.0 & Hz \\
Robot radius & 0.35 & m \\
Min obstacle distance & 0.4 & m \\
Preferred clearance & 1.0 & m \\
Free space weight & 0.3 & - \\
Obstacle penalty weight & 6.0 & - \\
Stuck movement threshold & 0.25 & m \\
Stuck radius threshold & 0.35 & m \\
PINN call timeout & 2.0 & s \\
\bottomrule
\end{tabular}
\end{table}

% TODO: Add Figure 3.4 - NSGA-II flowchart
% Show: Initialize → Evaluate → Non-dominated Sort → Crowding Distance →
%       Selection → Crossover → Mutation → Combine → Next generation
% Reference: Will be created as figures/nsga2_flowchart.pdf

% TODO: Add Figure 3.5 - Pareto front example
% 2D plot showing energy vs. path length for all solutions
% Color-code by front number (F1, F2, F3, ...)
% Reference: Will be created as figures/pareto_front_example.pdf


\section{UKF Localization}
\label{sec:meth_ukf}

The \acrfull{ukf} provides 5D state estimation by fusing wheel odometry, \acrshort{imu}, and magnetometer measurements. This section derives the complete \acrshort{ukf} formulation with enhancements for real-robot deployment.

% Code Reference: src/piec_ukf_localization/piec_ukf_localization/ukf_node.py

\subsection{State Space Model}

\paragraph{State Vector}
The robot state is represented by a 5-dimensional vector:
\begin{equation}
\vect{x} = \begin{bmatrix}
x \\ y \\ \theta \\ v \\ \omega
\end{bmatrix} \in \Real^5
\label{eq:state_vector}
\end{equation}

where $(x, y)$ is position, $\theta \in [-\pi, \pi]$ is orientation, $v$ is linear velocity, and $\omega$ is angular velocity.

\paragraph{Initial Covariance}
\begin{equation}
\mat{P}_0 = \text{diag}(0.2, 0.2, 0.05, 0.1, 0.01)
\label{eq:initial_covariance}
\end{equation}

\subsection{Motion Model}

The process model implements constant velocity dynamics with curved motion kinematics.

\paragraph{Straight Motion} ($|\omega| < 0.001$ rad/s):
\begin{equation}
\begin{aligned}
x_{k+1} &= x_k + v_k \cos(\theta_k) \Delta t \\
y_{k+1} &= y_k + v_k \sin(\theta_k) \Delta t \\
\theta_{k+1} &= \theta_k
\end{aligned}
\label{eq:motion_straight}
\end{equation}

\paragraph{Curved Motion} ($|\omega| \geq 0.001$ rad/s):
\begin{equation}
r = \frac{v_k}{\omega_k}
\label{eq:turn_radius}
\end{equation}

\begin{equation}
\begin{aligned}
x_{k+1} &= x_k + r \left[ \sin(\theta_k + \omega_k \Delta t) - \sin(\theta_k) \right] \\
y_{k+1} &= y_k + r \left[ \cos(\theta_k) - \cos(\theta_k + \omega_k \Delta t) \right] \\
\theta_{k+1} &= \theta_k + \omega_k \Delta t
\end{aligned}
\label{eq:motion_curved}
\end{equation}

\paragraph{Orientation Normalization}
To maintain $\theta \in [-\pi, \pi]$:
\begin{equation}
\theta_{\text{normalized}} = \text{atan2}(\sin(\theta), \cos(\theta))
\label{eq:angle_normalization}
\end{equation}

\paragraph{Velocity Dynamics}
\begin{equation}
\begin{aligned}
v_{k+1} &= v_k \\
\omega_{k+1} &= \omega_k
\end{aligned}
\label{eq:velocity_dynamics}
\end{equation}

Velocities are constant in the process model and updated through measurements.

\paragraph{Process Noise}
\begin{equation}
\vect{w}_k \sim \mathcal{N}(\vect{0}, \mat{Q})
\label{eq:process_noise}
\end{equation}

\begin{equation}
\mat{Q} = \text{diag}(0.1, 0.1, 0.05, 0.2, 0.2)
\label{eq:process_noise_cov}
\end{equation}

\textbf{Time step:} $\Delta t = 0.05$ s (20 Hz update rate).

\subsection{Measurement Models}

The system incorporates two measurement types:

\subsubsection{Pose Measurement (Wheel Odometry + IMU)}

\begin{equation}
\vect{z}_{\text{pose}} = \begin{bmatrix}
x \\ y \\ \theta
\end{bmatrix} + \vect{v}_{\text{pose}}
\label{eq:meas_pose}
\end{equation}

\textbf{Measurement function:}
\begin{equation}
h_{\text{pose}}(\vect{x}) = \begin{bmatrix}
x \\ y \\ \theta
\end{bmatrix}
\label{eq:h_pose}
\end{equation}

\textbf{Measurement noise covariance:}
\begin{equation}
\mat{R}_{\text{pose}} = \text{diag}(0.01, 0.01, 0.005)
\label{eq:R_pose}
\end{equation}

\subsubsection{Angular Velocity Measurement (IMU Gyroscope)}

\begin{equation}
z_{\omega} = \omega + v_{\omega}
\label{eq:meas_angular}
\end{equation}

\textbf{Measurement function:}
\begin{equation}
h_{\omega}(\vect{x}) = \omega
\label{eq:h_angular}
\end{equation}

\textbf{Measurement noise covariance:}
\begin{equation}
R_{\omega} = 0.01
\label{eq:R_angular}
\end{equation}

\subsection{Unscented Transform}

\subsubsection{Sigma Point Parameters}

\begin{equation}
\lambda = \alpha^2 (n + \kappa) - n
\label{eq:lambda}
\end{equation}

\begin{equation}
c = n + \lambda
\label{eq:c_param}
\end{equation}

where:
\begin{itemize}
    \item $\alpha = 10^{-3}$: Spread of sigma points
    \item $\beta = 2$: Optimal for Gaussian distributions
    \item $\kappa = 0$: Secondary scaling parameter
    \item $n = 5$: State dimension
\end{itemize}

\subsubsection{Sigma Point Generation}

Given current state $\vect{x}$ and covariance $\mat{P}$, generate $2n + 1 = 11$ sigma points:

\begin{equation}
\mathcal{X}_0 = \vect{x}
\label{eq:sigma_0}
\end{equation}

\begin{equation}
\mathcal{X}_i = \vect{x} + \left(\sqrt{c \mat{P}}\right)_{:,i}, \quad i = 1, \ldots, n
\label{eq:sigma_positive}
\end{equation}

\begin{equation}
\mathcal{X}_{i+n} = \vect{x} - \left(\sqrt{c \mat{P}}\right)_{:,i}, \quad i = 1, \ldots, n
\label{eq:sigma_negative}
\end{equation}

where $\sqrt{c \mat{P}}$ is the Cholesky decomposition: $\mat{S}\mat{S}\transpose = c\mat{P}$.

\subsubsection{Weights}

\paragraph{Mean Weights}
\begin{equation}
W_m^{(0)} = \frac{\lambda}{c}
\label{eq:weight_mean_0}
\end{equation}

\begin{equation}
W_m^{(i)} = \frac{1}{2c}, \quad i = 1, \ldots, 2n
\label{eq:weight_mean_i}
\end{equation}

\paragraph{Covariance Weights}
\begin{equation}
W_c^{(0)} = \frac{\lambda}{c} + (1 - \alpha^2 + \beta)
\label{eq:weight_cov_0}
\end{equation}

\begin{equation}
W_c^{(i)} = \frac{1}{2c}, \quad i = 1, \ldots, 2n
\label{eq:weight_cov_i}
\end{equation}

\subsection{Predict Step}

\paragraph{Propagate Sigma Points}
\begin{equation}
\mathcal{X}_{k+1|k}^{(i)} = f(\mathcal{X}_{k|k}^{(i)}, \Delta t), \quad i = 0, \ldots, 2n
\label{eq:propagate_sigma}
\end{equation}

where $f$ is the motion model (Equations~\ref{eq:motion_straight} or~\ref{eq:motion_curved}).

\paragraph{Predicted Mean}
\begin{equation}
\hat{\vect{x}}_{k+1|k} = \sum_{i=0}^{2n} W_m^{(i)} \mathcal{X}_{k+1|k}^{(i)}
\label{eq:predicted_mean}
\end{equation}

\paragraph{Predicted Covariance}
\begin{equation}
\mat{P}_{k+1|k} = \sum_{i=0}^{2n} W_c^{(i)} \left(\mathcal{X}_{k+1|k}^{(i)} - \hat{\vect{x}}_{k+1|k}\right)\left(\mathcal{X}_{k+1|k}^{(i)} - \hat{\vect{x}}_{k+1|k}\right)\transpose + \mat{Q}
\label{eq:predicted_covariance}
\end{equation}

\subsection{Update Step}

\paragraph{Measurement Sigma Points}
\begin{equation}
\mathcal{Z}_{k+1}^{(i)} = h(\mathcal{X}_{k+1|k}^{(i)}), \quad i = 0, \ldots, 2n
\label{eq:measurement_sigma}
\end{equation}

\paragraph{Predicted Measurement Mean}
\begin{equation}
\hat{\vect{z}}_{k+1} = \sum_{i=0}^{2n} W_m^{(i)} \mathcal{Z}_{k+1}^{(i)}
\label{eq:predicted_measurement}
\end{equation}

\paragraph{Innovation Covariance}
\begin{equation}
\mat{S}_{k+1} = \sum_{i=0}^{2n} W_c^{(i)} \left(\mathcal{Z}_{k+1}^{(i)} - \hat{\vect{z}}_{k+1}\right)\left(\mathcal{Z}_{k+1}^{(i)} - \hat{\vect{z}}_{k+1}\right)\transpose + \mat{R}
\label{eq:innovation_covariance}
\end{equation}

\paragraph{Cross-Covariance}
\begin{equation}
\mat{P}_{xz} = \sum_{i=0}^{2n} W_c^{(i)} \left(\mathcal{X}_{k+1|k}^{(i)} - \hat{\vect{x}}_{k+1|k}\right)\left(\mathcal{Z}_{k+1}^{(i)} - \hat{\vect{z}}_{k+1}\right)\transpose
\label{eq:cross_covariance}
\end{equation}

\paragraph{Kalman Gain}
\begin{equation}
\mat{K} = \mat{P}_{xz} \mat{S}_{k+1}^{-1}
\label{eq:kalman_gain}
\end{equation}

\paragraph{State Update}
\begin{equation}
\vect{x}_{k+1|k+1} = \hat{\vect{x}}_{k+1|k} + \mat{K} (\vect{z}_{k+1} - \hat{\vect{z}}_{k+1})
\label{eq:state_update}
\end{equation}

\paragraph{Covariance Update}
\begin{equation}
\mat{P}_{k+1|k+1} = \mat{P}_{k+1|k} - \mat{K} \mat{S}_{k+1} \mat{K}\transpose
\label{eq:covariance_update}
\end{equation}

\subsection{IMU Yaw Offset Correction}

Real \acrshort{imu} sensors often have systematic yaw biases and discontinuities at $\pm 180°$.

\paragraph{Quaternion to Yaw Conversion}
\begin{equation}
\begin{aligned}
\text{siny\_cosp} &= 2(q_w q_z + q_x q_y) \\
\text{cosy\_cosp} &= 1 - 2(q_y^2 + q_z^2) \\
\theta_{\text{raw}} &= \text{atan2}(\text{siny\_cosp}, \text{cosy\_cosp})
\end{aligned}
\label{eq:quaternion_to_yaw}
\end{equation}

\paragraph{Bias Removal and Offset Application}
\begin{equation}
\theta_{\text{corrected}} = \theta_{\text{raw}} - \theta_{\text{initial}} + \theta_{\text{offset}}
\label{eq:yaw_correction}
\end{equation}

\paragraph{Normalization}
\begin{equation}
\theta_{\text{final}} = \text{atan2}(\sin(\theta_{\text{corrected}}), \cos(\theta_{\text{corrected}}))
\label{eq:yaw_normalization}
\end{equation}

This ensures smooth transitions across the $\pm \pi$ boundary.

\subsection{Covariance Health Monitoring}

Long-duration missions require mechanisms to detect and recover from filter divergence.

\subsubsection{Health Check Criteria}

\paragraph{1. NaN/Inf Detection}
\begin{equation}
\text{Reset if} \quad \exists P_{ij} : P_{ij} \in \{\text{NaN}, \infty\}
\label{eq:health_nan}
\end{equation}

\paragraph{2. Trace Threshold}
\begin{equation}
\text{Reset if} \quad \text{trace}(\mat{P}) = \sum_{i=1}^{5} P_{ii} > 10.0
\label{eq:health_trace}
\end{equation}

Indicates excessive accumulated uncertainty.

\paragraph{3. Condition Number}
\begin{equation}
\kappa(\mat{P}) = \frac{\lambda_{\max}(\mat{P})}{\lambda_{\min}(\mat{P})} > 10^8 \implies \text{Reset}
\label{eq:health_condition}
\end{equation}

Indicates ill-conditioning (near-singularity).

\subsubsection{Regularization}

If minimum eigenvalue is too small:
\begin{equation}
\text{if} \quad \lambda_{\min}(\mat{P}) < 10^{-6}: \quad \mat{P} \leftarrow \mat{P} + (10^{-6} - \lambda_{\min} + 10^{-9}) \mat{I}
\label{eq:regularization}
\end{equation}

\subsubsection{Symmetry Enforcement}

Numerical errors can break symmetry:
\begin{equation}
\mat{P} \leftarrow \frac{1}{2}(\mat{P} + \mat{P}\transpose)
\label{eq:symmetry_enforcement}
\end{equation}

\subsubsection{Reset Covariance}

Upon divergence detection:
\begin{equation}
\mat{P}_{\text{reset}} = \text{diag}(0.01, 0.01, 0.01, 0.1, 0.1)
\label{eq:reset_covariance}
\end{equation}

State estimate is not reset; only covariance is reinitialized.

\subsection{Mahalanobis Distance Validation}

To reject outlier measurements:

\paragraph{Innovation}
\begin{equation}
\boldsymbol{\nu}_k = \vect{z}_k - \hat{\vect{z}}_k
\label{eq:innovation}
\end{equation}

\paragraph{Mahalanobis Distance}
\begin{equation}
d_M = \sqrt{\boldsymbol{\nu}_k\transpose \mat{S}_k^{-1} \boldsymbol{\nu}_k}
\label{eq:mahalanobis}
\end{equation}

\paragraph{Adaptive Threshold}
\begin{equation}
\tau_M = \begin{cases}
5.0 & \text{if rejection rate} \leq 0.3 \\
7.5 & \text{if rejection rate} > 0.3
\end{cases}
\label{eq:adaptive_threshold}
\end{equation}

\textbf{Rejection criterion:}
\begin{equation}
\text{Reject measurement if} \quad d_M > \tau_M
\label{eq:rejection_criterion}
\end{equation}

The adaptive threshold prevents excessive rejections during transient periods.

% TODO: Add Figure 3.6 - UKF sigma point distribution
% 2D plot showing state mean, covariance ellipse, and 11 sigma points
% Reference: Will be created as figures/ukf_sigma_points.pdf

% TODO: Add Figure 3.7 - Covariance trace over time
% Time series showing how trace(P) evolves, with reset events marked
% Reference: Will be created as figures/covariance_trace.pdf


\section{Adaptive DWA with Free-Space Awareness}
\label{sec:meth_dwa}

The local controller combines \acrfull{dwa} with free-space awareness, progressive rotation strategies, and anti-oscillation mechanisms for robust obstacle avoidance.

% Code Reference: src/piec_controller/piec_controller/
% - dwa_pinn_enhanced.py
% - controller_node.py
% - dynamic_dwa_complete.py

\subsection{DWA Fundamentals}

\acrshort{dwa} selects velocity commands by sampling the admissible velocity space and scoring forward-simulated trajectories.

\subsubsection{Velocity Space Constraints}

\paragraph{Kinematic Constraints}
\begin{equation}
\begin{aligned}
v &\in [0, v_{\max}] \\
\omega &\in [-\omega_{\max}, \omega_{\max}]
\end{aligned}
\label{eq:kinematic_constraints}
\end{equation}

\paragraph{Dynamic Constraints}
Considering acceleration limits over timestep $\Delta t$:
\begin{equation}
\begin{aligned}
v &\in [v_{\text{current}} - a_{\max} \Delta t, \, v_{\text{current}} + a_{\max} \Delta t] \\
\omega &\in [\omega_{\text{current}} - \alpha_{\max} \Delta t, \, \omega_{\text{current}} + \alpha_{\max} \Delta t]
\end{aligned}
\label{eq:dynamic_constraints}
\end{equation}

where $a_{\max} = 0.8$ m/s$^2$ and $\alpha_{\max} = 2.0$ rad/s$^2$.

\paragraph{Admissible Velocity Space}
\begin{equation}
\mathcal{V}_a = \{(v, \omega) : \text{kinematic} \cap \text{dynamic} \cap \text{collision-free}\}
\label{eq:admissible_space}
\end{equation}

\subsection{Trajectory Simulation}

For each velocity sample $(v, \omega) \in \mathcal{V}_a$, simulate forward trajectory over horizon $T_h = 2.0$ s:

\begin{equation}
\begin{aligned}
x(t + \Delta t) &= x(t) + v \cos(\theta(t)) \Delta t \\
y(t + \Delta t) &= y(t) + v \sin(\theta(t)) \Delta t \\
\theta(t + \Delta t) &= \theta(t) + \omega \Delta t
\end{aligned}
\label{eq:trajectory_simulation}
\end{equation}

Simulate at 10 Hz (10 steps over 2s horizon).

\subsection{Multi-Objective Trajectory Scoring}

Each trajectory is evaluated using five objectives:

\subsubsection{Objective 1: Goal Score}
\begin{equation}
O_{\text{goal}} = -\frac{d_{\text{final}}}{d_{\max} + 0.1}
\label{eq:dwa_goal}
\end{equation}

where $d_{\text{final}}$ is distance from trajectory endpoint to goal. Negative sign converts minimization to maximization.

\textbf{Weight:} $w_1 = 1.5$

\subsubsection{Objective 2: Speed Score}
\begin{equation}
O_{\text{speed}} = \frac{v}{v_{\max}}
\label{eq:dwa_speed}
\end{equation}

Rewards higher velocities for efficient navigation.

\textbf{Weight:} $w_2 = 0.7$

\subsubsection{Objective 3: Clearance Score}
\begin{equation}
O_{\text{clearance}} = \frac{d_{\min}}{d_{\text{safe}}}
\label{eq:dwa_clearance}
\end{equation}

where $d_{\min}$ is minimum distance to obstacles along trajectory and $d_{\text{safe}} = 0.8$ m.

\textbf{Weight:} $w_3 = 3.0$ (highest weight for safety)

\subsubsection{Objective 4: Path Alignment Score}
\begin{equation}
O_{\text{path}} = \frac{1}{N_{\text{traj}}} \sum_{i=1}^{N_{\text{traj}}} \frac{1}{1 + d_{\text{path}}(x_i, y_i)}
\label{eq:dwa_path}
\end{equation}

where $d_{\text{path}}(x, y)$ is distance from trajectory point to global path.

\textbf{Weight:} $w_4 = 1.0$

\subsubsection{Objective 5: Free-Space Score}
\begin{equation}
O_{\text{free}} = \frac{C_{\text{goal\_dir}}}{C_{\max}}
\label{eq:dwa_free}
\end{equation}

where $C_{\text{goal\_dir}}$ is clearance in goal direction and $C_{\max}$ is maximum clearance across all directions.

\textbf{Weight:} $w_5 = 0.8$

\subsubsection{Total DWA Score}
\begin{equation}
S_{\text{DWA}}(v, \omega) = \sum_{j=1}^{5} w_j \cdot O_j(v, \omega)
\label{eq:total_dwa_score}
\end{equation}

\textbf{Selection:}
\begin{equation}
(v^*, \omega^*) = \arg\max_{(v, \omega) \in \mathcal{V}_a} S_{\text{DWA}}(v, \omega)
\label{eq:velocity_selection}
\end{equation}

\subsection{Dynamic Velocity Limiting}

Standard \acrshort{dwa} uses fixed velocity limits. We implement obstacle-proximity-based dynamic limits.

\subsubsection{Speed Reduction Function}

\begin{equation}
v_{\text{adjusted}} = \begin{cases}
0 & d_{\text{obs}} \leq d_{\text{emergency}} \\
v \cdot \text{sigmoid}(d_{\text{obs}}, d_{\text{emergency}}, d_{\text{slow}}) & d_{\text{emergency}} < d_{\text{obs}} < d_{\text{slow}} \\
v \cdot (1 - 0.1 - 0.2\frac{d_{\text{safe}} - d_{\text{obs}}}{d_{\text{safe}}}) & d_{\text{slow}} \leq d_{\text{obs}} \leq d_{\text{safe}} \\
v & d_{\text{obs}} > d_{\text{safe}}
\end{cases}
\label{eq:speed_reduction}
\end{equation}

where:
\begin{itemize}
    \item $d_{\text{emergency}} = 0.2$ m: Emergency stop distance
    \item $d_{\text{slow}} = 0.6$ m: Begin slowing
    \item $d_{\text{safe}} = 0.8$ m: Safe distance threshold
\end{itemize}

\paragraph{Sigmoid Smoothing}
For smooth transitions:
\begin{equation}
\text{sigmoid}(d, d_1, d_2) = 0.3 + 0.7 \cdot \frac{1}{1 + e^{-8(t - 0.5)}}
\label{eq:sigmoid_smooth}
\end{equation}

where $t = \frac{d - d_1}{d_2 - d_1} \in [0, 1]$.

\subsection{Free-Space Awareness}

\subsubsection{Free-Space Direction Evaluation}

Sample 24 directions at 15° intervals:
\begin{equation}
\phi_i = \frac{2\pi i}{24}, \quad i = 0, 1, \ldots, 23
\label{eq:direction_samples}
\end{equation}

For each direction, compute clearance:
\begin{equation}
C(\phi_i) = \max\{r : \text{no obstacle at } (x + r\cos\phi_i, y + r\sin\phi_i), \, r \in [0, r_{\max}]\}
\label{eq:clearance_direction}
\end{equation}

\subsubsection{Free-Space Score}

For direction $\phi_i$:
\begin{equation}
F(\phi_i) = 0.8 \cdot C(\phi_i) + 0.2 \cdot \left(1 - \frac{|\phi_i - \phi_{\text{goal}}|}{\pi}\right)
\label{eq:free_space_direction_score}
\end{equation}

Combines clearance (80\%) with goal alignment (20\%).

\subsection{Progressive Rotation Thresholds}

To avoid excessive in-place rotation far from goal:

\begin{equation}
\tau_{\text{rotate}} = \begin{cases}
120° & d_{\text{goal}} > 2.0 \text{ m} \\
90° & 0.5 \text{ m} < d_{\text{goal}} \leq 2.0 \text{ m} \\
60° & d_{\text{goal}} \leq 0.5 \text{ m}
\end{cases}
\label{eq:progressive_thresholds}
\end{equation}

\textbf{Control logic:}
\begin{equation}
\text{Rotate in place} \iff |\theta_{\text{error}}| > \tau_{\text{rotate}}
\label{eq:rotation_condition}
\end{equation}

This prevents premature rotation when far from goal, allowing curved paths.

\subsection{Anti-Oscillation Mechanisms}

\subsubsection{Direction Locking}

Once a rotation direction is selected, maintain it for duration $T_{\text{lock}} = 2.0$ s:

\begin{equation}
\omega(t) = \begin{cases}
\text{sign}_{\text{locked}} \cdot \omega_{\max} & t - t_{\text{lock}} < T_{\text{lock}} \land |\theta_{\text{error}}| > \theta_{\min} \\
\text{sign}(\theta_{\text{error}}) \cdot \omega_{\max} & \text{otherwise}
\end{cases}
\label{eq:direction_locking}
\end{equation}

where $\theta_{\min} = 10°$ is minimum angle to maintain lock.

\subsubsection{Rotation Timeout}

If rotating in place exceeds timeout $T_{\text{timeout}} = 5.0$ s, force forward motion:

\begin{equation}
\text{if } t - t_{\text{rotate\_start}} > T_{\text{timeout}}: \quad \begin{cases}
v = \min(0.4 v_{\max}, 0.5 d_{\text{goal}}) \\
\omega = 0.5 \omega_{\max} \cdot \text{sign}(\theta_{\text{error}})
\end{cases}
\label{eq:rotation_timeout}
\end{equation}

This breaks deadlocks where pure rotation fails to resolve heading error.

\subsection{Heading Control Modes}

The controller operates in three modes depending on path geometry and heading error:

\subsubsection{Mode 1: Straight-Line Detection}

If path deviation is minimal:
\begin{equation}
|\Delta y| < 0.03 \text{ m} \land |\Delta x| > 0.2 \text{ m}
\label{eq:straight_line_condition}
\end{equation}

\textbf{Control:}
\begin{equation}
\begin{aligned}
v &= \min(0.7 v_{\max}, 0.5 d_{\text{goal}}) \\
\omega &= 0
\end{aligned}
\label{eq:straight_control}
\end{equation}

Forces zero angular velocity for efficient straight motion.

\subsubsection{Mode 2: Rotate-in-Place}

If heading error exceeds threshold:
\begin{equation}
|\theta_{\text{error}}| > \tau_{\text{rotate}}
\label{eq:rotate_condition}
\end{equation}

\textbf{Control:}
\begin{equation}
\begin{aligned}
v &= 0 \\
\omega &= \text{sign}(\theta_{\text{error}}) \cdot \omega_{\max}
\end{aligned}
\label{eq:rotate_control}
\end{equation}

with direction locking applied.

\subsubsection{Mode 3: Forward + Turn}

For moderate heading errors:
\begin{equation}
\frac{\tau_{\text{rotate}}}{2} < |\theta_{\text{error}}| \leq \tau_{\text{rotate}}
\label{eq:forward_turn_condition}
\end{equation}

\textbf{Control:}
\begin{equation}
\begin{aligned}
v &= v_{\text{base}} \cdot \alpha_{\text{align}} \\
\omega &= k_p \theta_{\text{error}}
\end{aligned}
\label{eq:forward_turn_control}
\end{equation}

where:
\begin{equation}
\alpha_{\text{align}} = \max\left(0.3, 1 - \frac{|\theta_{\text{error}}|}{\tau_{\text{rotate}}}\right)
\label{eq:alignment_factor}
\end{equation}

and $k_p = 1.5$ is proportional gain.

\subsection{PINN-Optimized Velocity Selection}

After \acrshort{dwa} selects velocity $(v, \omega)$, query \acrshort{pinn} for energy and stability:

\begin{algorithm}[h]
\caption{PINN Velocity Optimization}
\label{alg:pinn_velocity}
\begin{algorithmic}[1]
\State Generate trajectory proposal with $(v, \omega)$
\State Query PINN service: $(E, S) \leftarrow \text{PINN}(\text{trajectory})$
\If{$E > 50.0$ \textbf{or} $S < 0.7$}
    \State $v \leftarrow 0.8v$, $\omega \leftarrow 0.8\omega$ \Comment{Reduce speed}
\ElsIf{$E < 20.0$ \textbf{and} $S > 0.9$}
    \State $v \leftarrow \min(1.2v, v_{\max})$ \Comment{Increase speed}
\Else
    \State Keep $(v, \omega)$ unchanged
\EndIf
\State \textbf{return} $(v, \omega)$
\end{algorithmic}
\end{algorithm}

\textbf{Exception:} For straight paths ($|\Delta y| < 0.05$), bypass \acrshort{pinn} optimization to avoid unnecessary overhead.

\subsection{Algorithm Parameters}

\begin{table}[h]
\centering
\caption{DWA Controller Parameters}
\label{tab:dwa_params}
\begin{tabular}{lcc}
\toprule
\textbf{Parameter} & \textbf{Value} & \textbf{Unit} \\
\midrule
Max linear velocity & 1.0 & m/s \\
Min linear velocity & 0.05 & m/s \\
Max angular velocity & 0.7 & rad/s \\
Acceleration limit & 0.8 & m/s$^2$ \\
Deceleration limit & 1.2 & m/s$^2$ \\
Planning horizon & 2.0 & s \\
Velocity samples & 12 $\times$ 24 & - \\
Control frequency & 20 & Hz \\
Emergency stop distance & 0.2 & m \\
Slow down distance & 0.6 & m \\
Safe distance & 0.8 & m \\
Heading proportional gain & 1.5 & - \\
Max heading rate & 0.6 & rad/s \\
Rotate threshold (far) & 120 & ° \\
Rotate threshold (mid) & 90 & ° \\
Rotate threshold (close) & 60 & ° \\
Rotation timeout & 5.0 & s \\
Direction lock duration & 2.0 & s \\
Min angle for lock & 10 & ° \\
Free-space update rate & 3.0 & Hz \\
Free-space min clearance & 0.8 & m \\
\bottomrule
\end{tabular}
\end{table}

% TODO: Add Figure 3.8 - DWA velocity space sampling
% 2D plot showing (v, ω) sample grid with admissible/inadmissible regions
% Reference: Will be created as figures/dwa_velocity_space.pdf

% TODO: Add Figure 3.9 - Trajectory scoring visualization
% Show multiple candidate trajectories color-coded by score
% Reference: Will be created as figures/dwa_trajectories.pdf

\section{System Integration}
\label{sec:meth_integration}

This section describes how the five components integrate into a cohesive navigation system.

\subsection{ROS 2 Architecture}

The system is implemented as a collection of \acrshort{ros} 2 nodes communicating via topics and services.

\subsubsection{Node Structure}

\begin{enumerate}
    \item \textbf{ukf\_localization\_node}: Publishes filtered pose and covariance
    \item \textbf{path\_optimizer\_node}: Publishes global path (1 Hz)
    \item \textbf{pinn\_service\_node}: Provides energy/stability queries (ROS 2 service)
    \item \textbf{controller\_node}: Publishes velocity commands (20 Hz)
    \item \textbf{emergency\_stop\_node}: Safety monitor with override capability
\end{enumerate}

\subsubsection{Key Topics}

\begin{table}[h]
\centering
\caption{ROS 2 Topic Communication}
\label{tab:ros2_topics}
\begin{tabular}{lll}
\toprule
\textbf{Topic} & \textbf{Message Type} & \textbf{Rate} \\
\midrule
/odom & nav\_msgs/Odometry & 20 Hz \\
/imu/data & sensor\_msgs/Imu & 50 Hz \\
/scan & sensor\_msgs/LaserScan & 10 Hz \\
/ukf/pose & geometry\_msgs/PoseWithCovariance & 20 Hz \\
/global\_path & nav\_msgs/Path & 1 Hz \\
/cmd\_vel & geometry\_msgs/Twist & 20 Hz \\
/emergency\_stop & std\_msgs/Bool & 10 Hz \\
\bottomrule
\end{tabular}
\end{table}

\subsubsection{Service Interfaces}

\begin{itemize}
    \item \texttt{/pinn/evaluate\_trajectory}: Queries \acrshort{pinn} for energy and stability
    \begin{itemize}
        \item \textbf{Request:} Trajectory waypoints (positions, velocities, terrain data)
        \item \textbf{Response:} Energy (J/m), Stability (0-1)
        \item \textbf{Timeout:} 2.0 s
    \end{itemize}
\end{itemize}

\subsection{Control Loop Timing}

\paragraph{20 Hz Control Loop (50ms)}
\begin{enumerate}
    \item Receive UKF pose estimate
    \item Receive latest global path
    \item Query current laser scan
    \item Execute DWA trajectory evaluation
    \item (Optionally) Query PINN for energy optimization
    \item Publish velocity command
\end{enumerate}

\textbf{Typical breakdown:}
\begin{itemize}
    \item UKF update: 5ms
    \item DWA evaluation (288 trajectories): 30ms
    \item PINN query (if called): 5ms
    \item Path following logic: 5ms
    \item Safety checks: 2ms
    \item \textbf{Total:} $\sim$47ms (within 50ms budget)
\end{itemize}

\subsection{Failure Modes and Recovery}

\subsubsection{PINN Service Timeout}

If \acrshort{pinn} service exceeds 2.0s timeout:
\begin{enumerate}
    \item Use physics-based energy estimate: $E_{\text{fallback}} = 0.5v^2 + 0.2|\omega|$
    \item Log warning and continue with \acrshort{dwa}-only control
\end{enumerate}

\subsubsection{Localization Divergence}

If \texttt{trace}($\mat{P}$) $> 10.0$ or NaN detected:
\begin{enumerate}
    \item Reset covariance to $\mat{P}_{\text{reset}}$
    \item Increase process noise temporarily: $\mat{Q} \leftarrow 2\mat{Q}$ for 2 seconds
    \item Log divergence event
\end{enumerate}

\subsubsection{Path Optimizer Stuck Detection}

If stuck condition persists for 8 seconds:
\begin{enumerate}
    \item Increase exploration factor: $\alpha_{\text{explore}} = 0.6$
    \item Increase free-space weight: $w_{\text{free}} = 0.5$
    \item Execute escape maneuver (back up 0.5m, rotate to free-space direction)
    \item Replan with updated parameters
\end{enumerate}

\subsection{Parameter Configuration}

System parameters are stored in YAML configuration files:
\begin{itemize}
    \item \texttt{config/ukf\_params.yaml}: Covariance matrices, health thresholds
    \item \texttt{config/optimizer\_params.yaml}: NSGA-II population, objectives weights
    \item \texttt{config/dwa\_params.yaml}: Velocity limits, scoring weights
    \item \texttt{config/pinn\_params.yaml}: Model path, timeout, feature normalization
\end{itemize}

This allows tuning without recompilation.

% TODO: Add Figure 3.10 - System integration diagram
% Sequence diagram showing one complete control cycle with all nodes
% Reference: Will be created as figures/system_integration_sequence.pdf

\subsection{Summary}

This chapter presented the complete technical methodology of the \acrshort{piec} framework:
\begin{itemize}
    \item \textbf{PINN Surrogate Model:} Physics-informed energy prediction with 6-component decomposition
    \item \textbf{NSGA-II Optimizer:} Multi-objective path planning with 7 objectives
    \item \textbf{UKF Localization:} 5D state estimation with health monitoring and IMU calibration
    \item \textbf{Adaptive DWA:} Local control with free-space awareness and anti-oscillation
    \item \textbf{System Integration:} ROS 2 architecture with failure recovery
\end{itemize}

The next chapter discusses implementation details including simulation environment, real robot platform, and software architecture.

